\documentclass{article}
\usepackage{fancyhdr}
\usepackage[letterpaper, margin=1in]{geometry}
\usepackage[utf8]{inputenc}
\usepackage{amsmath}
\usepackage{circuitikz}
\title{4. Inductance and Capacitance \\  \large Introduction to Electric Circuits \\ Supplemental Instruction}
\author{Cooper McAllister}
\date{April 2 2026}
\begin{document}
\fancypagestyle{footerstyle} {
	\fancyhf{}
	\renewcommand{\headrulewidth}{0pt}
	\renewcommand{\footrulewidth}{0.4pt}
	\fancyfoot[L]{\textit{For personal study and students leading supplemental instruction, \\ with attribution given to the author. \\ Find more at coopermcallister.com}}
	\fancyfoot[R]{\thepage}
}
\pagestyle{footerstyle}
\maketitle
\thispagestyle{footerstyle}

\section{The Inductor}
The inductor is a circuit element indicated by either of the following symbols: \\ \\
\begin{circuitikz}[american, cute inductors] \draw
(0, 0) to[inductor] (3, 0)
(4, 0) to[american inductor] (7, 0)
;
\end{circuitikz} \\ \\
I will use the left symbol here.
\\ An inductor has an \textbf{inductance} L, measured in Henries (H). Expressed in SI units, the Henry is
$$H = \frac{V \cdot s}{A} = \frac{\Omega \cdot s}{rad} = \Omega \cdot s$$
\\ The equation relating voltage and current in an inductor is
$$v=L\frac{di}{dt}$$
Current cannot change instantly in an inductor. This is because for current to change instantly, its derivative would approach $\infty$, which would mean voltage would have to approach $\infty$, which is impossible.
\\ In steady-state (i.e. when nothing in the circuit is changing over time), an inductor behaves as a short. This is because the derivative of current is 0, so the voltage across the inductor is 0. \\ \\
\begin{circuitikz}[american, cute inductors] \draw
(0, 0) to[inductor, *-*] (3, 0)
(6, 0) to[short, *-*] (9, 0)
(4.2, -0.1) -- (4.8, -0.1) (4.2, 0.1) -- (4.8, 0.1) 
;
\end{circuitikz} $\qquad$ (steady state)
\\ Inductors can be combined in parallel and series using the same formulas as for resistors: \\ \\
\begin{circuitikz}[american, cute inductors] \draw
(0, 0) to[open] (0, 2) to[L=$L_1$] (3, 2) to[L=$L_2$] (5, 2) to[open] (7, 2) to[L=$L_n$] (9, 2)
-- (9, 0) -- (0, 0);
\node[draw=none] at (6, 2) {$\cdots$};
\end{circuitikz}
$\qquad \Longrightarrow \qquad$
\begin{circuitikz}[american, cute inductors] \draw
(0, 0) to[open] (0, 2)  to[L=$L_{eq}$] (3, 2) 
-- (3, 0) -- (0, 0)
;
\end{circuitikz} \\ \\
$$L_{eq} = L_1 + L_2 + \cdots + L_n$$
\\ \\
\begin{circuitikz}[american, cute inductors] \draw
(0, 0) to[open] (0, 2) to[short] (1, 2) to[L=$L_1$] (1, 0) (3, 2) to[L=$L_2$] (3, 0) (7, 2) to[L=$L_n$] (7, 0)
(7, 2) -- (1, 2) (7, 0) -- (0, 0);
\node[draw=none] at (5, 1) {$\cdots$};
\end{circuitikz}
$\qquad \Longrightarrow \qquad$
\begin{circuitikz}[american, cute inductors] \draw
(0, 0) to[open] (0, 2)  to[L=$L_{eq}$] (3, 2) 
-- (3, 0) -- (0, 0)
;
\end{circuitikz} \\ \\
$$L_{eq} = \frac{1}{\frac{1}{L_1} + \frac{1}{L_2} + \cdots + \frac{1}{L_n} }$$

\section{The Capacitor}
The capacitor is a circuit element indicated by either of the following symbols. The right symbol is used for polarized capacitors, which are not discussed in this course. \\ \\
\begin{circuitikz}[american, cute inductors] \draw
(0, 0) to[capacitor] (3, 0)
(4, 0) to[curved capacitor] (7, 0)
;
\end{circuitikz} \\ \\
A capacitor has a \textbf{capacitance} C, measured in Farads (F). Expressed in SI units, the Farad is
$$F = \frac{A \cdot s}{V} = \frac{s}{rad \cdot \Omega} = \frac{s}{\Omega}$$
The equation relating voltage and current in a capacitor is
$$i=C\frac{dv}{dt}$$
Voltage cannot change instantly in a capacitor. This is because for voltage to change instantly, its derivative would approach $\infty$, which would mean current would have to approach $\infty$, which is impossible.
\\ In steady-state, a capacitor behaves as an open circuit. This is because the derivative of voltage is 0, so the current through the inductor is 0. \\ \\
\begin{circuitikz}[american, cute inductors] \draw
(0, 0) to[capacitor, *-*] (3, 0)
(6, 0) to[open, *-*] (9, 0)
(4.2, -0.1) -- (4.8, -0.1) (4.2, 0.1) -- (4.8, 0.1) 
;
\end{circuitikz} $\qquad$ (steady state)
\\ Capacitors can be combined in parallel and series using the same formulas as for resistors, but using the opposite formulas for series and parallel: \\ \\
\begin{circuitikz}[american, cute inductors] \draw
(0, 0) to[open] (0, 2) to[C=$C_1$] (3, 2) to[C=$C_2$] (5, 2) to[open] (7, 2) to[C=$C_n$] (9, 2)
-- (9, 0) -- (0, 0);
\node[draw=none] at (6, 2) {$\cdots$};
\end{circuitikz}
$\qquad \Longrightarrow \qquad$
\begin{circuitikz}[american, cute inductors] \draw
(0, 0) to[open] (0, 2)  to[C=$C_{eq}$] (3, 2) 
-- (3, 0) -- (0, 0)
;
\end{circuitikz} \\ \\
$$C_{eq} = \frac{1}{\frac{1}{C_1} + \frac{1}{C_2} + \cdots + \frac{1}{C_n} }$$
\\ \\
\begin{circuitikz}[american, cute inductors] \draw
(0, 0) to[open] (0, 2) to[short] (1, 2) to[C=$C_1$] (1, 0) (3, 2) to[C=$C_2$] (3, 0) (7, 2) to[C=$C_n$] (7, 0)
(7, 2) -- (1, 2) (7, 0) -- (0, 0);
\node[draw=none] at (5, 1) {$\cdots$};
\end{circuitikz}
$\qquad \Longrightarrow \qquad$
\begin{circuitikz}[american, cute inductors] \draw
(0, 0) to[open] (0, 2)  to[C=$C_{eq}$] (3, 2) 
-- (3, 0) -- (0, 0)
;
\end{circuitikz} \\ \\
$$C_{eq} = C_1 + C_2 + \cdots + C_n$$


\section{First-Order Time Response}
Circuits with DC sources and \textit{only one} capacitor or inductor exhibit first-order response. If a change takes place at time $t=0$ (such as a switch opening or closing), any current or voltage in the circuit, $m(t)$, can be found as:
$$m(t) = (m(0^+) - m(\infty))e^{\frac{-t}{\tau}} + m(\infty)$$
where
\\$m(0^+)$ is the quantity evaluated immediately after the change takes place,
\\$m(\infty)$ is the quantity evaluated a very long time after the change takes place, and
\\$\tau$ is the time constant:
\\
\\For a circuit with a capacitor, $$\tau = R_{th}C$$
\\For a circuit with an inductor, $$\tau = \frac{L}{R_{th}}$$
\\$R_{th}$ is the thevenin resistance of the circuit as seen from the capacitor or inductor.
\\Note that $t=0^-$ is used to denote the time immediately before the change takes place (while the circuit is at DC), $t=0^+$ is used to denote the time immediately after the change takes place (the start of the response), and $t=\infty$ is used to denote the time a very long time after the change takes place (when the circuit is once again at DC).
\\A general procedure for finding the response of capacitor voltage or inductor current in a first-order circuit is as follows:
\begin{enumerate}
\item Draw the circuit before the change takes place (at $t=0^-$) and replace capacitors and inductors with their DC equivalent.
\item Find:
\begin{itemize}
\item For a circuit with a capacitor: $V_C(0^-)$, the voltage across the capacitor. Then, $V_C(0^-) = V_C(0^+)$ (because capacitor voltage cannot change instantly).
\item For a circuit with an inductor: $I_L(0^-)$, the current through the inductor. Then, $I_L(0^-) = I_L(0^+)$ (because inductor current cannot change instantly).
\item Note that $V_C(0^-)$ or $I_L(0^-)$ may be given if there is an initial capacitor voltage or inductor current.
\end{itemize}
\item Draw the circuit after the change takes place, replace capacitors and inductors with their DC equivalent, and find $V_C(\infty)$ or $I_L(\infty)$.
\item Find the time constant using the cirucit after the change takes place.
\item Substitute $m(0^+)$, $m(\infty)$, and $\tau$ with the proper values in the above equation to obtain the time response.
\end{enumerate}

\newpage
\section*{References}
\begin{itemize}
\item M. Iordache. Class Lectures, Topic: ``Electric circuits.'' EEGR 2053, School of Engineering and Engineering Technology, LeTourneau University, 2024.
\item M. Iordache. 2024. EEGR 2053 Electric circuits [PowerPoint slides]. Available: \\ https://mviordache.name/EEGR2053/index.html
\item T. L. Floyd and D. M. Buchla, Principles of Electric Circuits: Conventional Current, 10th ed. Hoboken, NJ: Pearson Education, 2020.
\item O. Ortiz. 2025. ENGR 2053 Intro to electric circuits [PowerPoint slides].
\end{itemize}



\end{document}